% ===========================================================================
\section{Discussion}
\label{sec:discussion}
% ===========================================================================

\subsection{One sector, two emission accounts}

The central result is that this sector is \SIrange{5.4}{6.5}{\percent} of Pakistan's
greenhouse gases across the two dung configurations while carrying a configuration-
dependent but consistently leading share of its primary particulate, and that both
readings of the sector are correct.

It is tempting to resolve that tension by arguing that one accounting is wrong---that biomass
\COtwo{} should count, or that particulate should be given a \COe{} equivalent. Neither move is
defensible. The biomass memo convention is an accounting rule that prevents double
counting: biogenic \COtwo{} is held outside Energy because any net change in biomass carbon
stocks is accounted in AFOLU---whether Pakistan's fuelwood harvest is in fact balanced by
regrowth is a land-sector question this paper does not evaluate. Black carbon's metrics genuinely are contested, and converting it
would imply a comparability with long-lived gases that the science does not support.

The tension is not an artefact to be resolved but a fact about the sector: it burns a lot
of carbon whose net atmospheric effect is booked to the land sector, and emits a lot of
primary particulate. An inventory that reports only one of those describes half of it.

\subsection{What the closure buys, and what it does not}

The dung-cake transfer removes a specific failure. Before it, \SI{1.259}{\mega\tonne} of
excreted nitrogen a year left agriculture's manure balance correctly and entered nothing---an
orphaned flow that both sectors' totals were individually consistent with and that was wrong
between them.

What it does not buy is accuracy in the dung quantity itself---or even certainty that
households are the destination. The transfer guarantees that the diverted mass is
accounted for exactly once, as household combustion or as a typed unresolved remainder;
it does not establish that the \SI{28}{\percent} and
\SI{30}{\percent} burned-fuel shares are right, nor that everything they divert is burned
in dwellings. The HIES demand-side reading, a factor of \num{5.4} lower, is why the
reference books only what households report burning, why this paper
reports the all-diverted configuration beside it
(Table~\ref{tab:configs}), and why it restates the identity as
$D_{\text{diverted}} = D_{\text{residential}} + D_{\text{other/unresolved}}$. The closure
makes the two sectors share one assumption about what was diverted instead of holding two
inconsistent ones---an
improvement in consistency, while the unresolved remainder carries the open evidence
honestly.

That distinction generalises to the other closures in this inventory
\citep{ahmad2026national}: each removes an inconsistency and none removes an unknown.

\subsection{The firewood decision}

The firewood question used to be whether to include a figure carrying its source's own
estimate flag. It is now a different and better question: how far to trust a national
household survey.

The answer this sector gives is: as the primary quantity, at \SI{35}{\percent}. The HIES
total is a design-weighted measurement reproduced within \SI{6}{\percent} by an independent
round six years earlier, and the supply-side estimate it replaces lands within
\SI{7}{\percent} of it---three numbers with different biases agreeing at the ten-percent
level, which is more corroboration than most activity data in this inventory can claim. The
residual \SI{35}{\percent} prices the systematic risks a consumption survey cannot see:
recall over a one-month window, out-of-dwelling use, enterprise wood booked to households.

The honest reading of the reported figure is a range with two axes, and the
configuration axis is the wider one: \SI{24.2}{} (the reference) to
\SI{29.5}{\mega\tonne} (configuration~L), not a band
around the reference alone. Within the reference, the firewood axis spans \SI{20.8}{} to
\SI{25.4}{}, with the lower bound being the no-firewood case---which is not ``the
metered fuels alone'': it still carries the non-market dung stream, and the metered
fossil fuels by themselves support about \SI{18.4}{\mega\tonne}.

\subsection{Implications}

\textbf{Household solid fuel is the largest covered source of primary particulate.} Under
the reference and configuration L, residential combustion emits
\numrange{49}{109} times the primary \PM{} attributed to road exhaust. Ranking interventions
requires exposure, cost and abatement-effectiveness evidence that this inventory does not
provide.

\textbf{The two national surveys this sector needed both exist, and both are now held.}
HIES quantities replaced the FAOSTAT firewood estimate and supply the reference dung
quantity---a factor of five below what the nitrogen ledger diverts, with the difference
carried as unresolved rather than reconciled away. And the Multi-Tier Framework stove typology
\citep{worldbank2024mtf} supplies the national stove distribution: \SI{44.3}{\percent} of households cook mainly on clean-fuel
stoves, \SI{71.8}{\percent} of rural households on traditional or three-stone/open-fire
stoves, with stacking at \SI{6.2}{\percent}. The stove-mix sensitivity now prices that
evidence; the remaining need is the fuel-by-device-by-end-use crosswalk the report does
not publish, which is what separates this inventory from technology resolution.

\textbf{Fuel switching moves the two inventories in opposite proportions.} Replacing biomass
with piped gas would \emph{increase} this sector's reported greenhouse gases, because gas
\COtwo{} is reportable and biomass \COtwo{} is not, while sharply reducing its particulate. Any
policy evaluation that scores such a switch on the national inventory alone will score it as a
loss. That is a genuine accounting artefact rather than a physical result, and it is the
clearest case in this inventory for reporting both layers together.
